\documentclass[twocolumn]{article}
\usepackage{index-flmrecomposed.87f71783}
\begin{document}
\title{A document of several files}
\author{Ann Author}
\date{30 September 2026}
\maketitle
\flmHeading[toc=true]
\section{Introduction}{\flmLabels{\label{ref1}}}

This paragraph begins in the primary file, continues with words of an included file, which names \flmUDproject{} through the house
package's macro,
 and ends in the primary
file, as one paragraph. The method is in \flmRef{ref2}.\flmXNfootnote{A note of the primary
file.}

\flmHeading[toc=true]
\section{The \emph{Method}}{\flmLabels{\label{ref2}}}
A section kept in a file of its own, back to \flmRef{ref1}.\flmXNfootnote{A note of the
included file.}

\flmHeading[toc=true]
\subsection{Details}{\flmLabels{}}

The details, read relative to the file that names them:
\begin{itemize}
\item one;
\item two.
\end{itemize}



\flmHeading[toc=true]
\section{Conclusion}{\flmLabels{}}

The primary file goes on after the section kept in a file of its own, and its notes are
counted after that file's.\flmXNfootnote{A note of the primary file after the included file.}
%%
\end{document}
